\documentclass[twocolumn,10pt]{article}
%\usepackage[utf8]{utf8}
\usepackage[pt-BR]{babel}
\usepackage{amsmath,amssymb,amsfonts}
\usepackage{graphicx}
\usepackage{cite}
\usepackage{hyperref}
\usepackage{listings}
\usepackage{xcolor}
\usepackage{geometry}

\geometry{
    a4paper,
    top=20mm,
    bottom=20mm,
    left=18mm,
    right=18mm
}

\definecolor{codegreen}{rgb}{0,0.6,0}
\definecolor{codegray}{rgb}{0.5,0.5,0.5}
\definecolor{codepurple}{rgb}{0.58,0,0.82}
\definecolor{backcolour}{rgb}{0.95,0.95,0.92}

\lstdefinestyle{matlabstyle}{
    backgroundcolor=\color{backcolour},   
    commentstyle=\color{codegreen},
    keywordstyle=\color{blue},
    numberstyle=\tiny\color{codegray},
    stringstyle=\color{codepurple},
    basicstyle=\ttfamily\footnotesize,
    breakatwhitespace=false,         
    breaklines=true,                 
    captionpos=b,                    
    keepspaces=true,                 
    numbers=left,                    
    numbersep=5pt,                  
    showspaces=false,                
    showstringspaces=false,
    showtabs=false,                  
    tabsize=2
}

\lstset{style=matlabstyle}

\title{\textbf{Relatório de Iniciação Científica -- Identificação Conversor Buck}}
\author{\textbf{Luiz Felipe Souza de Lima Silva}\\
Engenharia Elétrica -- UFPA\\
\textit{Assunto: Conversor Buck}}
\date{10 de Maio de 2026}

\begin{document}

\maketitle

\begin{abstract}
Neste relatório foi disposta toda a simulação e coleta de dados do conversor BUCK, com simulação realizada no software Simulink.
\\[0.5em]
\textbf{Palavras-chave:} Buck, Controle, Modelagem, Eletrônica de Potência, Identificação, Simulação, Simulink, Conversor Buck.
\end{abstract}

\section{Introdução}
A relação entre Machine Learning (ML) e Identificação de Sistemas (SI) evoluiu de campos distintos para uma convergência tecnológica profunda, onde as técnicas de aprendizagem de máquina são agora ferramentas naturais para a modelação de sistemas dinâmicos complexos \cite{pillonetto2023deep}.

\section{Fundamentação Teórica}
O procedimento foi composto por uma coleta de dados e posteriormente uma abordagem para criação de um modelo a partir dos dados obtidos.

\subsection{Simulação}
A simulação foi realizada no software Simulink, onde todas as ferramentas (toolboxes) foram obtidas através da subpasta de aplicação denominada Power Simscape.

O conversor Buck foi montado em sua forma clássica.

\subsection{Coleta de Dados}
A coleta de dados foi realizada extraindo um conjunto de dados de entrada e saída do sistema real, importando-os para um arquivo ``arquivo.csv'' e posteriormente adicionando-os ao MATLAB.

São elas:
\begin{itemize}
    \item \textbf{Saída} ($y(k)$): Tensão dissipada na carga resistiva.
    \item \textbf{Entrada} ($u(k)$): Sinal PRBS aplicado no gate da chave MOSFET.
\end{itemize}

\subsection{Modelo}
A partir da coleta de dados, obtivemos um modelo do sistema partindo da ideia de modelagem por identificação, também conhecida como modelagem em caixa preta, onde o modelo é obtido através de um conjunto de dados de entrada e saída \cite{aguirre2015introducao}.

\subsubsection{A Escolha do Modelo}
O modelo ARX (do inglês \textit{AutoRegressive with eXogenous input} -- AutoRegressivo com entrada eXógena) é considerado a estrutura fundamental para a identificação de sistemas lineares. Ele descreve a relação entre o sinal de entrada $u(k)$ e a saída $y(k)$ de um sistema através de uma equação de diferenças lineares.

\subsubsection{Estrutura}
A equação de diferenças do modelo ARX pode ser definida por:
\begin{equation}
y(k) = \sum_{i=1}^{n} a_i y(k-i) + \sum_{j=1}^{n} b_j u(k-j)
\end{equation}

Para o modelo obtido através do conjunto de dados, considerando dois coeficientes de entrada e saída, teremos:
\begin{equation}
y(k) = a_1 y(k-1) + a_2 y(k-2) + b_1 u(k-1) + b_2 u(k-2)
\end{equation}

\subsection{Matriz dos Regressores}
Temos que, para $k \ge 3$:
\begin{align*}
y(3) &= a_1 y(2) + a_2 y(1) + b_1 u(2) + b_2 u(1) \\
y(4) &= a_1 y(3) + a_2 y(2) + b_1 u(3) + b_2 u(2) \\
&\;\;\vdots \\
y(k) &= a_1 y(k-1) + a_2 y(k-2) + b_1 u(k-1) + b_2 u(k-2)
\end{align*}

Por definição, a partir do sistema linear, podemos gerar a seguinte representação matricial $A \cdot X = Y$:

\begin{equation}
\underbrace{
\begin{bmatrix} 
y(3) \\ 
y(4) \\ 
\vdots \\ 
y(N) 
\end{bmatrix}
}_{Y}
=
\underbrace{
\begin{bmatrix} 
y(2) & y(1) & u(2) & u(1) \\
y(3) & y(2) & u(3) & u(2) \\
\vdots & \vdots & \vdots & \vdots \\
y(N-1) & y(N-2) & u(N-1) & u(N-2)
\end{bmatrix}
}_{A}
\underbrace{
\begin{bmatrix} 
a_1 \\ 
a_2 \\ 
b_1 \\ 
b_2 
\end{bmatrix}
}_{X}
\end{equation}

Em que $A$ é denotada como a matriz dos regressores.

Isolando $X$ através do método dos mínimos quadrados ($A \cdot X = Y$), obtemos:
\begin{equation}
X = (A^T \cdot A)^{-1} \cdot A^T \cdot Y
\end{equation}

Com o algoritmo dos mínimos quadrados executado, basta encontrar os coeficientes da matriz $X$.

\subsection{Sinal de Entrada $u(k)$}
Para o sinal de entrada, é aplicado um ruído branco, gerando um sinal PRBS de entrada.

\subsection{Sinal de Saída $y(k)$}
O sinal de saída é coletado como a resposta do modelo. No caso do conversor Buck, a resposta é o nível de tensão na carga.

Portanto, obtém-se a função de transferência discreta:
\begin{equation}
G(q) = \frac{b_1 q^{-1} + b_2 q^{-2}}{1 - a_1 q^{-1} - a_2 q^{-2}}
\end{equation}

O modelo discreto obtido pode ser comparado com um modelo ideal.

\section{Resultados e Discussão}

\subsection{Espectros dos Sinais}
Para analisar a composição em frequência dos sinais utilizados, foi gerado o espectro de magnitude do sinal de entrada PRBS e do sinal de saída. A Figura \ref{fig:espectros} apresenta o resultado da Transformada de Fourier aplicada aos sinais, permitindo observar as características em frequência do sistema.

\begin{figure}[htbp]
    \centering
    \includegraphics[width=\linewidth]{IdentREAL_01.png}
    \caption{Espectros de magnitude do sinal de entrada e de saída.}
    \label{fig:espectros}
\end{figure}

\subsection{Relação Entrada-Saída}
O processo de identificação utilizou um sinal PRBS para excitar de forma diversificada as dinâmicas de operação do conversor Buck. A Figura \ref{fig:saidaentrada} ilustra a relação temporal dos dados coletados, com a entrada apresentando pulsos ruidosos ($u$) e a resposta real refletindo a variação correspondente do nível de tensão de saída ($v$).

\begin{figure}[htbp]
    \centering
    \includegraphics[width=\linewidth]{IdentREAL_03.png}
    \caption{Ensaio de Identificação do Conversor Buck: Sinal de entrada PRBS (superior) e Saída Real (inferior).}
    \label{fig:saidaentrada}
\end{figure}

\subsection{Saída Real vs Saída ARX}
Com base nos dados coletados, estimou-se a matriz de parâmetros do modelo ARX via Mínimos Quadrados. A Figura \ref{fig:comparacao} comprova graficamente o grau de precisão do modelo paramétrico frente ao comportamento da planta. O sinal simulado (ARX) acompanha quase perfeitamente as variações apresentadas no sistema real validado.

\begin{figure}[htbp]
    \centering
    \includegraphics[width=\linewidth]{IdentREAL_05.png}
    \caption{Gráfico de comparação temporal: Sinal do modelo ARX vs. Saída do sistema real.}
    \label{fig:comparacao}
\end{figure}

\subsection{Funções de Transferência ($G_z$ e $G_{ze}$) e Ajuste ($R^2$)}
Após o processamento algorítmico num tempo de amostragem de $50\,\mu\text{s}$, extraiu-se a função de transferência de tempo discreto idealizada ($G_z$) baseada em modelagem matemática e a nova função de transferência discreta estimada ($G_{ze}$) gerada pelo algoritmo, dadas por:

\begin{equation}
G_z = \frac{1.785z + 1.785}{z^2 - 1.768z + 0.8869}
\end{equation}

\begin{equation}
G_{ze} = \frac{0.01594z - 0.03983}{z^2 - 1.729z + 0.7917}
\end{equation}

Para verificar analiticamente a coerência de aproximação desse modelo linear de caixa-preta ($G_{ze}$), obteve-se um Coeficiente de Determinação ($R^2$) equivalente a:
\begin{equation}
R^2 = 0.9741
\end{equation}

Esse valor comprova matematicamente a fidelidade visual observada na Figura \ref{fig:comparacao}, indicando que o modelo ARX extraído consegue representar aproximadamente $97.41\%$ de todas as variações contidas na saída real da planta física original após o alinhamento dos dados.

\subsection{Análise de Resíduos}
A análise residual do sistema confirma a precisão do modelo ARX gerado. O erro de predição ($e = y - y_p$) possui uma média de $-4.9658 \times 10^{-4}$. O comportamento temporal desse erro é exposto na Figura \ref{fig:residuo}.

\begin{figure}[htbp]
    \centering
    \includegraphics[width=\linewidth]{IdentREAL_07.png}
    \caption{Resíduos / Erro de Predição $e(t)$ ao longo das amostras.}
    \label{fig:residuo}
\end{figure}

\subsection{Mapeamento de Polos e Zeros}
O mapa de polos e zeros do sistema discreto também foi levantado e comparado entre a função de transferência analítica ($G_z$) e a estimada ($G_{ze}$). 
\begin{itemize}
    \item O modelo idealizado $G_z$ possui polos complexos conjugados em $0.8839 \pm 0.3249i$ e um zero em $-1$.
    \item O modelo estimado $G_{ze}$ possui polos em $0.8646 \pm 0.2100i$ e um zero em $2.4997$.
\end{itemize}
A Figura \ref{fig:poloszeros} expõe graficamente essa disposição.

\begin{figure}[htbp]
    \centering
    \includegraphics[width=\linewidth]{IdentREAL_09.png}
    \caption{Mapa de polos e zeros (Pole-Zero Map) comparativo dos modelos discretos.}
    \label{fig:poloszeros}
\end{figure}

\section{Conclusão}
\begin{center}
    \textit{falta escrever}
\end{center}

\begin{thebibliography}{9}

\bibitem{pillonetto2023deep}
G.~Pillonetto, A.~Aravkin, D.~Gedon, L.~Ljung, A.~H.~Ribeiro, e T.~B.~Schön,
``Deep networks for system identification: A survey,''
\textit{Automatica}, 2023.

\bibitem{aguirre2015introducao}
L.~A.~Aguirre,
\textit{Introdução à Identificação de Sistemas: Técnicas Lineares e Não-Lineares Aplicadas a Sistemas Reais},
4ª~ed.\ Belo Horizonte: Editora UFMG, 2015.

\end{thebibliography}

\end{document}